% ===== Джерела до розділу 17. Ключі з префіксом ch17: =====
% Позиції ядра (bib-core.tex) тут НЕ дублюються.

% --- Управління та зрілість ---
\bibitem{ch17:nist-csf} National Institute of Standards and Technology. The NIST Cybersecurity Framework (CSF) 2.0. NIST CSWP~29. Gaithersburg: NIST, 2024. DOI: 10.6028/NIST.CSWP.29.

\bibitem{ch17:iso27001} International Organization for Standardization. ISO/IEC~27001: Information Security Management Systems --- Requirements. Geneva: ISO, 2022.

\bibitem{ch17:c2m2} U.S. Department of Energy. Cybersecurity Capability Maturity Model (C2M2), Version~2.1. Washington: DOE, 2022.

% --- Економіка та управління ризиками ---
\bibitem{ch17:anderson} Anderson R., Moore T. The Economics of Information Security. \emph{Science}. 2006. Vol.~314, №~5799. P.~610--613. DOI: 10.1126/science.1130992.

\bibitem{ch17:gordon-loeb} Gordon L. A., Loeb M. P. The Economics of Information Security Investment. \emph{ACM Transactions on Information and System Security}. 2002. Vol.~5, №~4. P.~438--457.

\bibitem{ch17:fair} Freund J., Jones J. Measuring and Managing Information Risk: A FAIR Approach. Waltham: Butterworth-Heinemann, 2014. 408~p.

% --- Люди та культура ---
\bibitem{ch17:awareness} Bada M., Sasse A. M., Nurse J. R. C. Cyber Security Awareness Campaigns: Why Do They Fail to Change Behaviour? \emph{International Conference on Cyber Security for Sustainable Society}. 2015. P.~118--131.

\bibitem{ch17:verizon} Verizon. Data Breach Investigations Report (DBIR). Verizon Business, 2024. (щорічний аналітичний звіт про чинник людини в інцидентах).

% --- Стратегія та стійкість ---
\bibitem{ch17:resilience} Ross R., Pillitteri V., Graubart R., Bodeau D., McQuaid R. Developing Cyber-Resilient Systems: A Systems Security Engineering Approach. NIST Special Publication 800-160, Vol.~2, Rev.~1. Gaithersburg: NIST, 2021.

% --- Українські джерела ---
\bibitem{ch17:strategy-ua} Про рішення Ради національної безпеки і оборони України «Про Стратегію кібербезпеки України». Указ Президента України. Київ, 2021.

\bibitem{ch17:governance-ua} Організаційно-управлінські засади побудови системи кіберзахисту підприємства в умовах сучасних загроз. \emph{Кібербезпека: освіта, наука, техніка}. 2024. URL: \url{https://csecurity.kubg.edu.ua/index.php/journal/article/view/830}.
